In this chapter we will verify approaches described in Chapter \ref{ch:implementation}
and perform evaluation of their computation time.

To verify approaches and validate the results of reachability and configuration checking 
 and to demonstrate the feasibility of 
the counterexample-guided control approach, we complement the formal analysis 
with the simulation. 
We will use the simulation platform described in Chapter \ref{ch:simulation}.
By comparing the outcomes of formal verification with simulation 
results, we aim to assess both the soundness and the practical utility of 
model-checking-based methods for modular self-replicating robots.



\section{Validation of the counterexample based robot control}
\label{ch:validationCounterexample}

To validate the counterexample based approach, introduced in 
chapter \ref{ch:Counterexample} we will integrate the counterexample generation 
pipeline into ROS2-Gazebo simulation environment.
We choose robot blocks configuration described in Chapter \ref{ch:robotSimulation}, 
robot initial position and target point with coordinates $[50,50,10]$. 
These initial parameters reflected in the lines 1-3 of the
 Listing \ref{lst:validCounterInit} lines.

\begin{lstlisting}[language=PythonMyStyle, caption={Counterexample class application example}, 
    label={lst:validCounterInit}]
angle_vars = ["block1.yaw", "block2.pitch", "block3.pitch", "block4.pitch", "block5.yaw"]
sequence = ["base", "yaw", "pitch", "pitch", "pitch", "yaw"]
lengths = [10, 20, 30, 30, 10, 20]
robot = Robot(sequence, lengths)
generator = RobotNusmvGenerator(robot, output_smv_path="./smv/")
generator.set_target_point(50, 50, 10)
generator.generate_from_template("project_robot")
result = Counterexample.run_nusmv_file("./smv/robot.smv")
counterexample = Counterexample(result, len(sequence))
angles = counterexample.get_angles(angle_vars)
counterexample.plot_trace(-1, show_plot=True)
\end{lstlisting}
Trace derived from the counterexample consist out of 9 steps and 
shown in Figure \ref{fig:validCounterexamplePlot}
\begin{figure}[h!]
    \centering
    \includegraphics[height=10cm]{validCounterexamplePlot.png} 
    \caption{Trace derived from counterexample of reachability of point.}
    \label{fig:validCounterexamplePlot}
\end{figure}

Resulting joint angles derived from the trace presented as a list steps,
where each step is a list of control angles for each joint:
\[ [0, -10, -10, -10, -10] \]
\[ [0, -20, -20, -20, -20] \]
\[ [0, -30, -30, -30, -30] \]
\[ [0, -30, -40, -40, -40] \]
\[ [0, -30, -50, -50, -50] \]
\[ [0, -30, -60, -50, -60] \]
\[ [0, -30, -70, -50, -60] \]
\[ [0, -30, -80, -50, -60] \]

These angles were published as a sequence of control messages 
for the robot's joints using the rotate\_joints() function implemented in 
the RobotController class, which is described in Section \ref{ch:robotSimulation}.
This function ensures that each joint of the simulated robot receives a  
command corresponding to the angles obtained from the verification results,
thereby reproducing in simulation the trajectory.

The result of this is shown in 
Figure \ref{fig:validCounterSimulation}, which illustrates the robot successfully performing the 
movement sequence derived from the counterexample.

\begin{figure}[h!]
    \centering
    % \includegraphics[height=4cm]{counterSimulation/0.png} 
    \includegraphics[height=4.5cm]{counterSimulation/1.png} 
    \includegraphics[height=4.5cm]{counterSimulation/2.png} 
    \includegraphics[height=4.5cm]{counterSimulation/3.png} 
    \includegraphics[height=4.5cm]{counterSimulation/4.png} 
    \includegraphics[height=4.5cm]{counterSimulation/5.png} 
    \includegraphics[height=4.5cm]{counterSimulation/6.png} 
    \includegraphics[height=4.5cm]{counterSimulation/7.png} 
    \includegraphics[height=4.5cm]{counterSimulation/8.png} 
    \includegraphics[height=4.5cm]{counterSimulation/9.png} 
    \caption{Simulation steps to achieve target point from counterexample}
    \label{fig:validCounterSimulation}
\end{figure}

Before launching the simulation in Gazebo, it was necessary to align the 
coordinate systems between the counterexample data and the simulation environment.
The coordinate system orientation used in the NuSMV model differs from that of 
Gazebo. Therefore, a preliminary transformation 
step was applied to adjust the  rotation of the robot and 
the target. This alignment 
was critical for correctly mapping the counterexample trajectory to 
the simulated robot's physical motion. Once this transformation was 
established, the robot was able to execute the motion plan as described by 
the verified specification. 

For the target point considered in this experiment, the resulting distance 
between the robot's end effector and the target position after motion completion 
was approximately 0.16 simulation units. This residual distance represents 
the cumulative error introduced by several factors. 
% The total discrepancy between 
% the target and the actual position of the end effector is primarily composed 
% of  the discretization error in the angular representation within NuSMV 
% and the geometric offsets present in the Gazebo simulation model, 
% particularly related to the simplifications made in the mechanical description 
% of the robot's modular blocks. Across a series of ten experiments with 
% different reachable target points, the maximum observed positional deviation 
% was about 27\% of one voxel unit (the basic unit of spatial resolution in 
% the simulation). Despite this discrepancy, the results consistently confirmed 
% that the trajectories derived from counterexamples could be translated into 
% feasible robot motions.
%

Total discrepancy between the target point and the center of the 
end effector consists of the angle discretization error in NuSMV and 
the geometric offset of the blocks in the simulation model. 
Maximal discrepancy during simulation of ten different reachable points 
was 27\% of length unit in the simulation (size of the voxel).

The simulation process also revealed several challenges and 
limitations in the current implementation of the counterexample-
based control mechanism. 
The main observed issues are summarized below:

\begin{enumerate}
    \item \textbf{Non-orthogonality of the end effector during mechanical attachment.}\\
    During attachment actions, the end effector did not remain 
    orthogonal to the side surfaces of the parts being assembled. 
    This misalignment  reduce the effectiveness of physical attachment or 
    cause mechanical stress on the joints.
    
    \item \textbf{Potential self-intersection of robot segments during movement.}\\
    While following the motion trace generated in NuSMV, the robot occasionally 
    approached poses that would lead to self-intersection in 
    the physical model. 
    This issue arises because our model used in NuSMV lacks explicit 
    spatial collision constraints.
    
    \item \textbf{Positional error due to angle discretization.}\\
    The angular variables in the NuSMV model were represented with a 
    discrete resolution (e.g., in 10$^\circ$ or smaller increments). 
    This quantization introduces rounding errors, which accumulate along 
    multi-joint kinematic chains and lead to noticeable deviations in end-effector positioning.
    
    \item \textbf{Geometric inaccuracy due to simplified structural modeling.}\\
    The NuSMV model employs a simplified geometric abstraction of the robot
    structure, which does not fully capture the precise spatial offsets 
    between certain block types.
    Particularly, in our model blocks with a rotational axis, offset along the longitudinal direction. 
\end{enumerate}

To overcome these limitations and improve the accuracy of counterexample-
based control, several refinements and future system enhancements
are proposed:

\begin{enumerate}
    \item The NuSMV model can be extended with additional constraints that 
    enforce orthogonal alignment of the terminal blocks of the robot relative 
    to the surfaces of the parts being attached. 
    This adjustment will ensure that the resulting trajectories maintain proper 
    contact geometry during assembly.
    
    \item By incorporating geometric consistency checks or collision-avoidance predicates 
        into the model generation process, 
    it is possible to prevent self-intersecting configurations from being 
    produced during verification. 
    Such constraints could be encoded as temporal logic formulas that prohibit 
    conflicting joint combinations.
    
    \item The angular discretization can be progressively refined in an iterative process. 
    The model-checking trace from a coarse discretization can be used as a 
    baseline, followed by repeated verification steps with smaller angle 
    increments. 
    This refinement strategy will balance computational complexity and accuracy.
    
    \item 
    The robot physical structure in the NuSMV representation should be further refined 
    to reflect the exact spatial arrangement of the blocks, 
    including any offsets or rotations in the mechanical joints. 
    This would reduce systematic discrepancies between the verified and simulated trajectories.
\end{enumerate}

In summary, the integration of counterexample-based motion generation with 
Gazebo simulation demonstrates the feasibility of using formal verification outputs 
to control physical or simulated self-replicating robots. 
Although the current implementation still exhibits measurable deviations due 
to model simplifications and discretization, 
the results confirm that verified behaviors can be meaningfully reproduced 
in a physical simulation environment. 
Continued refinement of the model fidelity and constraint formulation will 
further improve the correspondence between the symbolic verification domain 
and the physical reality of robotic assembly.



% The simulation also revealed the following challenges in the current implementation 
% of counterexample-based control:
%
% \begin{enumerate}
%     \item non-orthogonality of the end effector to the sides of the part during mechanical attachment.
%     \item potential self-intersection of the robot during movement along the path
%     of the robot in NSMV for orthogonal connection to parts.
%     \item position error due to angle discretization.
%     \item error due to simplified representation of the robot structure in NYSMV.
% \end{enumerate}
%
% These issues can be resolved in the following system improvements by:
%
% \begin{enumerate}
% \item introducing additional constraints on the blocks close to the end 
%     of the robot in NYSMV for orthogonal connection to the parts.
% \item introducing constraints on self-intersections of robot blocks during movement.
% \item iterative reduction of the discretization angle with the introduction of 
%     restrictions as the trace is refined.
% \item refinement of the robot's geometry in the model.
% \end{enumerate}


\section{Validation of the robot configuration checking}
\label{ch:validationConfiguration}

In order to prove consistence of configurations checking workflow described 
in Chapter \ref{ch:configurationChecking} we will first compare model checking results for 
both sequential checking with python and aggregated checking with NuSMV. 
And then  we will simulate one of reachable configurations in Gazebo, using simulation 
platform described in Chapter \ref{ch:simulation}  to  validate 
the results and prove that the target point is reachable for chosen configuration.


\subsection{Comparison of individual and aggregated model checking}

We will perform sequential configuration model checking for robots which consists of 
block of types: "yaw", "pitch", 
for target point $[20,20,20]$ and target region:

$x=\{-10,10\}, y=\{-10,10\}, z=\{0,10\}$

We will use \texttt{configuration\_checking} class, described in Chapter \ref{ch:configurationCheckingPy}.
and the working script which is located in \texttt{configuration\_checking/}
\texttt{validation\_individual\_configurations.py} script.

For robot consisting of  3 blocks using individual checking 
we  have no reachable configuration, 
the list of all possible configurations are given in Table \ref{tab:conf3individ}.

% Checking reachability of point: ['20', '20', '20']:
% 1/4:base_pitch_pitch_:  not reachable   [0.027696 sec]
% 2/4:base_pitch_yaw_:    not reachable   [0.011898 sec]
% 3/4:base_yaw_pitch_:    not reachable   [0.016229 sec]
% 4/4:base_yaw_yaw_:      not reachable   [0.010167 sec]

\begin{table}[h!]
\begin{tabular}{|l|l|l|l|l|l|}
\hline
N  & Configration                           & Point  & Time,sec   & Region  & Time,sec    \\ \hline
1 & base\_pitch\_pitch &  not reachable   & 0.0276    & not reachable  & 0.0117 \\ \hline
2 & base\_pitch\_yaw   &  not reachable   & 0.0118    & not reachable  & 0.0144 \\ \hline
3 & base\_yaw\_pitch   &  not reachable   & 0.0162    & not reachable  & 0.0232 \\ \hline
4 & base\_yaw\_yaw     &  not reachable   & 0.0101    & not reachable  & 0.0131 \\ \hline
\end{tabular}
\caption{Configurations reachability checking for robot size 3 for target point and target region.}
\label{tab:conf3individ}
\end{table}

For the aggregated model checking with NuSMV model adjusted for robot length 3 we  
also had no reachable configurations, and the computation time was 0.047455 sec.
This result corresponds to individual checking outcome.

For robot consisting of 4 blocks, using individual checking 
we get five reachable configuration,
the list of all possible configurations are given in Table \ref{tab:conf4individ}.

% 1/8:base_pitch_pitch_pitch_:     reachable      [0.126555 sec]
% 2/8:base_pitch_pitch_yaw_:       reachable      [0.035818 sec]
% 3/8:base_pitch_yaw_pitch_:       reachable      [0.187626 sec]
% 4/8:base_pitch_yaw_yaw_:        not reachable   [0.0123 sec]
% 5/8:base_yaw_pitch_pitch_:       reachable      [0.107405 sec]
% 6/8:base_yaw_pitch_yaw_:         reachable      [0.024083 sec]
% 7/8:base_yaw_yaw_pitch_:        not reachable   [0.062486 sec]
% 8/8:base_yaw_yaw_yaw_:  not reachable   [0.01128 sec]

\begin{table}[h!]
\begin{tabular}{|l|l|l|l|l|l|}
\hline
N  & Configration                           & Point  & Time,sec   & Region  & Time,sec    \\ \hline
1 & base\_pitch\_pitch\_pitch & reachable      & 0.1265   &  not reachable   & 0.3181  \\ \hline
2 & base\_pitch\_pitch\_yaw   & reachable      & 0.0358   &  not reachable   & 0.0402  \\ \hline
3 & base\_pitch\_yaw\_pitch   & reachable      & 0.1876   &  not reachable   & 0.2137  \\ \hline
4 & base\_pitch\_yaw\_yaw     & not reachable  & 0.0123   &  not reachable   & 0.0164  \\ \hline
5 & base\_yaw\_pitch\_pitch   & reachable      & 0.1074   &  not reachable   & 0.1253  \\ \hline
6 & base\_yaw\_pitch\_yaw     & reachable      & 0.0240   &  not reachable   & 0.0235  \\ \hline
7 & base\_yaw\_yaw\_pitch     & not reachable  & 0.0624   &  not reachable   & 0.0654  \\ \hline
8 & base\_yaw\_yaw\_yaw       & not reachable  & 0.0112   &  not reachable   & 0.0171  \\ \hline
\end{tabular}
\caption{Configurations reachability checking for robot size 4 for target point and target region.}
\label{tab:conf4individ}
\end{table}

And for the aggregated model checking with NuSMV model adjusted for robot length=4 we will 
have a following result for the reachability of point $[20,20,20]$ :

['base', 'pitch', 'pitch', 'yaw']
Time: 0.699504 sec

This result corresponds to the robot configuration number 2 in Table \ref{tab:conf4individ}
for individual checking.


% self.length=5
% 1/16:base_pitch_pitch_pitch_pitch_:      reachable      [1.637597 sec]
% 2/16:base_pitch_pitch_pitch_yaw_:        reachable      [0.274722 sec]
% 3/16:base_pitch_pitch_yaw_pitch_:        reachable      [2.074039 sec]
% 4/16:base_pitch_pitch_yaw_yaw_:  reachable      [0.044345 sec]
% 5/16:base_pitch_yaw_pitch_pitch_:        reachable      [2.124321 sec]
% 6/16:base_pitch_yaw_pitch_yaw_: not reachable   [0.209023 sec]
% 7/16:base_pitch_yaw_yaw_pitch_:  reachable      [1.063883 sec]
% 8/16:base_pitch_yaw_yaw_yaw_:   not reachable   [0.012656 sec]
% 9/16:base_yaw_pitch_pitch_pitch_:        reachable      [0.987572 sec]
% 10/16:base_yaw_pitch_pitch_yaw_:         reachable      [0.224217 sec]
% 11/16:base_yaw_pitch_yaw_pitch_:         reachable      [2.172816 sec]
% 12/16:base_yaw_pitch_yaw_yaw_:  not reachable   [0.022455 sec]
% 13/16:base_yaw_yaw_pitch_pitch_:        not reachable   [0.455085 sec]
% 14/16:base_yaw_yaw_pitch_yaw_:  not reachable   [0.065636 sec]
% 15/16:base_yaw_yaw_yaw_pitch_:  not reachable   [0.122896 sec]
% 16/16:base_yaw_yaw_yaw_yaw_:    not reachable   [0.01179 sec]
%

Results for model checking of robots consisting of 5 blocks,  
shown in Table \ref{tab:conf5individ}.

\begin{table}[h!]
\begin{tabular}{|l|l|l|l|l|l|}
\hline
N  & Configration                           & Point  & Time,sec   & Region  & Time,sec    \\ \hline
1  & base\_pitch\_pitch\_pitch\_pitch  & reachable        & 1.6375   & not reachable    & 7.0634  \\ \hline
2  &  base\_pitch\_pitch\_pitch\_yaw   & reachable        & 0.2747   & not reachable    & 0.3352  \\ \hline
3  &  base\_pitch\_pitch\_yaw\_pitch   & reachable        & 2.0740   & not reachable    & 3.4262  \\ \hline
4  &  base\_pitch\_pitch\_yaw\_yaw     & reachable        & 0.0443   & not reachable    & 0.0442  \\ \hline
5  &  base\_pitch\_yaw\_pitch\_pitch   & reachable        & 2.1243   & reachable        & 8.6275  \\ \hline
6  &  base\_pitch\_yaw\_pitch\_yaw     & not reachable    & 0.2090   & not reachable    & 0.2357  \\ \hline
7  &  base\_pitch\_yaw\_yaw\_pitch     & reachable        & 1.0638   & not reachable    & 1.1828  \\ \hline
8  &  base\_pitch\_yaw\_yaw\_yaw       & not reachable    & 0.0126   & not reachable    & 0.0183  \\ \hline
9  &  base\_yaw\_pitch\_pitch\_pitch   & reachable        & 0.9875   & reachable        & 1.3594  \\ \hline
10 &  base\_yaw\_pitch\_pitch\_yaw     & reachable        & 0.2242   & reachable        & 0.2629  \\ \hline
11 &  \cellcolor{green!20}base\_yaw\_pitch\_yaw\_pitch     & reachable        & 2.1728   & reachable        & 4.3590  \\ \hline
12 &  base\_yaw\_pitch\_yaw\_yaw       & not reachable    & 0.0224   & not reachable    & 0.0462  \\ \hline
13 &  base\_yaw\_yaw\_pitch\_pitch     & not reachable    & 0.4550   &  not reachable   & 0.5033  \\ \hline
14 &  base\_yaw\_yaw\_pitch\_yaw       & not reachable    & 0.0656   & not reachable    & 0.0753  \\ \hline
15 &  base\_yaw\_yaw\_yaw\_pitch       & not reachable    & 0.1228   & not reachable    & 0.1382  \\ \hline
16 &  base\_yaw\_yaw\_yaw\_yaw         & not reachable    & 0.0117   & not reachable    & 0.0177  \\ \hline
\end{tabular}
\caption{Configurations reachability checking for robot size 5 for target point and target region.}
\label{tab:conf5individ}
\end{table}


For the aggregated model checking with NuSMV model adjusted for robot length 5 we will 
have a following result for the reachability of target point:

['base', 'pitch', 'yaw', 'pitch', 'pitch']

Calculation time was 22.020454 sec

All aggregated model checking results corresponds to the individual checking, 
that proves the consistency of the developed approach.
% Rerion {-10,10} {-10,10} {0,10}

% \begin{table}[h!]
% \begin{tabular}{|l|l|l|l|l|l|}
% \hline
% N  & Configration                           & Point  & Time,sec   & Region  & Time,sec    \\ \hline
% 1 & base\_pitch\_pitch\_pitch\_pitch\_pitch &  reachable     & 169.71c & not reachable     & 1007.4  \\ \hline
% 2 & base\_pitch\_pitch\_pitch\_pitch\_yaw &   reachable      & 1.8283  & not reachable     & 7.3235  \\ \hline
% 3 & base\_pitch\_pitch\_pitch\_yaw\_pitch &   reachable      & 70.105  & reachable      & 466.35  \\ \hline
% 4 & base\_pitch\_pitch\_pitch\_yaw\_yaw &    not reachable     & 0.2860  & not reachable      & 0.3562  \\ \hline
% 5 & base\_pitch\_pitch\_yaw\_pitch\_pitch &   reachable      & 50.922  & reachable      & 226.78  \\ \hline
% 6 & base\_pitch\_pitch\_yaw\_pitch\_yaw &     reachable      & 2.4502  & reachable      & 10.687  \\ \hline
% 7 & base\_pitch\_pitch\_yaw\_yaw\_pitch &     reachable      & 17.172  & not reachable     & 38.146  \\ \hline
% 8 & base\_pitch\_pitch\_yaw\_yaw\_yaw &      not reachable     & 0.0463  & not reachable      & 0.0862  \\ \hline
% 9 & base\_pitch\_yaw\_pitch\_pitch\_pitch &   reachable      & 55.968  & reachable      & 207.57  \\ \hline
% 10 & base\_pitch\_yaw\_pitch\_pitch\_yaw &    reachable      & 6.5920  & reachable      & 10.631  \\ \hline
% 11 & base\_pitch\_yaw\_pitch\_yaw\_pitch &    reachable      & 147.78  & reachable      & 889.28  \\ \hline
% 12 & base\_pitch\_yaw\_pitch\_yaw\_yaw &      reachable      & 0.2445  & not reachable     & 0.2833  \\ \hline
% 13 & base\_pitch\_yaw\_yaw\_pitch\_pitch &    reachable      & 18.397  & not reachable     & 29.298  \\ \hline
% 14 & base\_pitch\_yaw\_yaw\_pitch\_yaw &      reachable      & 1.1950  & not reachable     & 1.6801  \\ \hline
% 15 & base\_pitch\_yaw\_yaw\_yaw\_pitch &     not reachable     & 2.7777  & not reachable      & 3.8429  \\ \hline
% 16 & base\_pitch\_yaw\_yaw\_yaw\_yaw &       not reachable     & 0.0128  & not reachable      & 0.0335  \\ \hline
% 17 & base\_yaw\_pitch\_pitch\_pitch\_pitch &  reachable      & 25.389  & reachable      & 38.161  \\ \hline
% 18 & base\_yaw\_pitch\_pitch\_pitch\_yaw &    reachable      & 2.2987  & reachable      & 3.1978  \\ \hline
% 19 & base\_yaw\_pitch\_pitch\_yaw\_pitch &    reachable      & 46.454  & reachable      & 238.72  \\ \hline
% 20 & base\_yaw\_pitch\_pitch\_yaw\_yaw &     not reachable     & 0.2329  & not reachable      & 0.2708  \\ \hline
% 21 & base\_yaw\_pitch\_yaw\_pitch\_pitch &    reachable      & 45.252  & reachable      & 150.38  \\ \hline
% 22 & base\_yaw\_pitch\_yaw\_pitch\_yaw &      reachable      & 5.0098  & reachable      & 7.8665  \\ \hline
% 23 & base\_yaw\_pitch\_yaw\_yaw\_pitch &     not reachable     & 17.253  & reachable       & 27.916  \\ \hline
% 24 & base\_yaw\_pitch\_yaw\_yaw\_yaw &       not reachable     & 0.0482  & not reachable      & 0.0510  \\ \hline
% 25 & base\_yaw\_yaw\_pitch\_pitch\_pitch &    reachable      & 2.8130  & not reachable     & 4.0847  \\ \hline
% 26 & base\_yaw\_yaw\_pitch\_pitch\_yaw &      reachable      & 0.5298  & not reachable     & 0.5949  \\ \hline
% 27 & base\_yaw\_yaw\_pitch\_yaw\_pitch &      reachable      & 7.4332  & not reachable     & 27.393  \\ \hline
% 28 & base\_yaw\_yaw\_pitch\_yaw\_yaw &       not reachable     & 0.0784  & not reachable      & 0.0887  \\ \hline
% 29 & base\_yaw\_yaw\_yaw\_pitch\_pitch &     not reachable     & 1.1831  & not reachable      & 1.3677  \\ \hline
% 30 & base\_yaw\_yaw\_yaw\_pitch\_yaw &       not reachable     & 0.1470  & not reachable      & 0.3001  \\ \hline
% 31 & base\_yaw\_yaw\_yaw\_yaw\_pitch &       not reachable     & 0.6042  & not reachable      & 0.6404  \\ \hline
% 32 & base\_yaw\_yaw\_yaw\_yaw\_yaw &         not reachable     & 0.0119  & not reachable      & 0.0188  \\ \hline
% \end{tabular}
% \caption{Configurations reachability checking for robot size 6.}
% \label{tab:conf5individ}
% \end{table}
% Time: 1057.268079 sec

The important observation is that in configurations where the first 
block after base is pitch block, the configuration
is equivalent to a base-yaw-pitch..., as if we would have an additional yaw block after base,
this property defined by base block structure,
and reflects the arbitrary initial rotation of a robot.


\subsection{Simulation of eligible configurations}

% We create a new list of reachable configurations for robots of length 5 and target point 
% with coordinates[20,20,20]. Table \ref{tab:con5target20}


% 6  &  base\_pitch\_yaw\_pitch\_yaw   & \cellcolor{green!20}reachable        & \cellcolor{green!20}0.2200     \\ \hline
% \begin{table}[h!]
%     \begin{center}
%         
% \begin{tabular}{|l|l|l|l|}
% \hline
% N  & Configration                           & Point  & Time,sec   \\ \hline
% 1  & base\_pitch\_pitch\_pitch\_pitch  & reachable        & 1.9822         \\ \hline
% 2  &  base\_pitch\_pitch\_pitch\_yaw   & reachable        & 0.2775         \\ \hline
% 3  &  base\_pitch\_pitch\_yaw\_pitch   & reachable        & 2.2516         \\ \hline
% 4  &  base\_pitch\_pitch\_yaw\_yaw     & not reachable    & 0.0418         \\ \hline
% 5  &  base\_pitch\_yaw\_pitch\_pitch   & reachable        & 2.4341         \\ \hline
% 6  &  base\_pitch\_yaw\_pitch\_yaw     & reachable        & 0.2340          \\ \hline
% 7  &  base\_pitch\_yaw\_yaw\_pitch     & not reachable    & 1.0739         \\ \hline
% 8  &  base\_pitch\_yaw\_yaw\_yaw       & not reachable    & 0.0161         \\ \hline
% 9  &  base\_yaw\_pitch\_pitch\_pitch   & reachable        & 1.2558         \\ \hline
% 10 &  base\_yaw\_pitch\_pitch\_yaw     & reachable        & 0.2558         \\ \hline
% 11 &  \cellcolor{green!20}base\_yaw\_pitch\_yaw\_pitch     & reachable        & 2.7031         \\ \hline
% 12 &  base\_yaw\_pitch\_yaw\_yaw       & not reachable    & 0.0246         \\ \hline
% 13 &  base\_yaw\_yaw\_pitch\_pitch     & not reachable    & 0.4640         \\ \hline
% 14 &  base\_yaw\_yaw\_pitch\_yaw       & not reachable    & 0.0673         \\ \hline
% 15 &  base\_yaw\_yaw\_yaw\_pitch       & not reachable    & 0.1293         \\ \hline
% 16 &  base\_yaw\_yaw\_yaw\_yaw         & not reachable    & 0.0133         \\ \hline
% \end{tabular}
% \end{center}
% \caption{Configurations reachability checking for robot size 5.}
% \label{tab:con5target20}
% \end{table}


From the list of reachable configurations presented in Table~\ref{tab:conf5individ}, 
one configuration was selected for validation using the simulation platform. 
The chosen configuration, highlighted in green in the table, corresponds to 
the following sequence of blocks:
\texttt{[base, yaw, pitch, yaw, pitch]}. 
A robot with this structure is capable of reaching the specified target 
point according to the model-checking results.

To implement this configuration in the simulation, a dedicated 
xacro file was created, and located in 
\texttt{/urdf/robot\_yaw\_pitch\_yaw\_pitch.xacro}, in \cite{GitHub}. 
The primary modification in this file concerns the block sequence of the robot, 
as shown in Listing~\ref{lst:configCheckingValidationRobot}. 
This modification ensures that the simulated robot matches
the verified structural configuration used in NuSMV.

\begin{lstlisting}[language=NuSMV, caption={Structure of reachable configuration of a robot}, 
    label={lst:configCheckingValidationRobot}]
<xacro:voxel name="block0" origin="0 0 0" color="${red}"/>
<xacro:block_ver block_num="1" par_num="0"  x="0" y="0" z="0" />
<xacro:block_fix block_num="2" par_num="1"  x="0" y="0" z="1" />
<xacro:block_hor block_num="3" par_num="2"  x="0" y="1" z="1" />
<xacro:block_ver block_num="4" par_num="3"  x="0" y="0" z="-0.5" />
<xacro:block_fix block_num="6" par_num="4"  x="0" y="0" z="1" />
<xacro:block_hor block_num="7" par_num="6"  x="0" y="-1" z="1" />
<xacro:block_fix block_num="8" par_num="7"  x="0" y="0" z="0.5" />
<xacro:block_end block_num="11" par_num="8" x="0" y="0" z="0.5" />
\end{lstlisting}

A counterexample trace for this configuration and the corresponding target points 
was obtained using the script \texttt{counterexample\_yaw\_pitch\_yaw\_pitch\_robot.py} 
available in~\cite{GitHub}. 
The resulting joint angles extracted from the counterexample are:
\texttt{[10, -70, -80, -70]}. 
These values represent the rotation commands for each joint in degrees and were later applied in the simulation to reproduce the verified trajectory.

For the purpose of simulation and validation, a dedicated ROS~2 launch file named 
\texttt{robot\_config\_validation.launch.py} was developed. 
This file initializes the Gazebo environment, loads the robot description, 
and executes the counterexample-based control sequence to move the robot toward the target point.

The results of the simulation are presented in Figure~\ref{fig:configCheckingValidationRobot}. 
The figure illustrates the initial configuration of the robot and its final pose after performing 
the movement sequence derived from the counterexample trace. 
As can be observed, the robot successfully reached the target point, 
confirming the feasibility of the configuration identified through model checking.


% From a list of reachable configurations in Table \ref{tab:conf5individ}
% we will choose one to validate the result using simulation platform. 
% The chosen configuration highlighted in green.
% As an example we will choose following configuration of blocks:
% [base, yaw, pitch, yaw, pitch]. 
% Robot with such configuration can reach a target point.
%
%  First of all we will create a xacro file for robot configuration, this file presented 
%  in /urdf/robot\_yaw\_pitch\_yaw\_pitch.xacro in \cite{GitHub}. The main changes are in 
%  block sequence, shown in Listing \ref{lst:configCheckingValidationRobot}.
%

% We recieved a counterexample for this configuration and target points
% using script counterexample\_yaw\_pitch\_yaw\_pitch\_robot.py in \cite{GitHub}.
% The resulting angles for joints are: [10, -70, -80, -70].
%
% In order to check this configuration of a robot we created a separate ROS2 launch file:
% robot\_config\_validation.launch.py.
%
% The simulation result of the counterexample based movement to a target point 
% configuration shown in Figure \ref{fig:configCheckingValidationRobot}

\begin{figure}[h!]
    \centering
    \includegraphics[height=9cm]{config_validation1.png} 
    \includegraphics[height=9cm]{config_validation2.png} 
    \caption{Validation of reachable configuration of a robot, initial and target pose}
    \label{fig:configCheckingValidationRobot}
\end{figure}

% Thus, in this section, we have shown that the results of individual and 
% aggregated verification of NuSMV models for all possible robot configurations up 
% to length 5 are consistent with each other. We also provided an example 
% of a simulation of one of the obtained models, which confirmed the results. 
%

In summary, this section demonstrates that the outcomes of both individual 
and aggregated verification 
of NuSMV models for all robot configurations up to a length of five blocks 
are consistent with each other. 
Furthermore, the provided simulation example validates one of the verified 
configurations, 
confirming that the model-checking results accurately correspond to 
physically realizable robot movements within the simulation environment.

\subsection{Comparison of computational time for model checking}
\label{ch:timeAnalysis}

To analyze the computational complexity of configuration checking and to compare 
the aggregated and sequential verification approaches, a series of 
experiments were conducted. 
In the first part of the experiment, a single target point with coordinates 
$x = 20$, $y = 20$, $z = 20$ was used for reachability verification. 
The results of these experiments were presented in Tables~\ref{tab:conf3individ}--\ref{tab:conf5individ}.

In the second part, the total model-checking time was evaluated for 
a target region defined by the coordinate ranges 
$x \in \{-10, 10\}$, $y \in \{-10, 10\}$, and $z \in \{0, 10\}$. 
The variation of the total verification time for point reachability with respect to 
the number of blocks in the robot 
is shown in Figure~\ref{fig:ConfigCheckingPointTimeSize}. 
% In order to analyze the complexity of the configuration checking and compare 
% aggregated checking approach with sequential we will perform a number of experiments.
%
% At first part used a  point with coordinates: x=20, y=20, z=20 as a target.
%
% Results are presented in Tables \ref{tab:conf3individ}-\ref{tab:conf5individ}. 
%
% The changing of total time required for model checking with  number of blocks in 
% robot shown in the Figure \ref{fig:ConfigCheckingPointTimeSize}

\begin{figure}[h!]
    \centering
    \includegraphics[height=9cm]{Time_size_plot_point.png} 
    \caption{Total time of model checking with number of blocks in robot for target point.}
    \label{fig:ConfigCheckingPointTimeSize}
\end{figure}

Similarly, the dependence of the total verification time for region reachability 
on the robot size is illustrated in Figure~\ref{fig:ConfigCheckingRegionTimeSize}.
% The second part used a total time of model checking of a region:
% $x=\{-10,10\}, y=\{-10,10\}, z=\{0,10\}]$
% The changing of total time required for model checking with  number of blocks in 
% robot shown in the Figure \ref{fig:ConfigCheckingPointTimeSize}

\begin{figure}[h!]
    \centering
    \includegraphics[height=9cm]{Time_size_plot_region.png} 
    \caption{Total time of model checking with number of blocks in robot for  target region.}
    \label{fig:ConfigCheckingRegionTimeSize}
\end{figure}


% Sequential checking of the configuration has the following features:
% \begin{enumerate}
%     \item It gives a full list of all reachable configurations
%     \item It allows to further analyze the generated configuration
%     \item It is slower in general than aggregated approach
%  \end{enumerate}
%
% Aggregated configurations checking from the other hand:
% \begin{enumerate}
%     \item Is faster than sequential checking for common case
%     \item Gives only the first encountered reachable configuration, 
%         can find others only after negating the found one.
%     \item Allows constraint specific blocks by number in sequence 
%         with type defining variable in NuSMV
%  \end{enumerate}
In both graphs, the model checking time grows exponentially with robot 
size. This is a typical behavior for model checking problems, as the 
state space explosion becomes dominant when system complexity increases. 
For smaller robot sizes (3–4 modules), the total checking time is nearly 
negligible in both point and region cases. 
However, as the robot size reaches 5 and especially 6 modules, the computation 
time rises  in the point-target case up to approximately 700
 seconds, and in the region-target case exceeding around 2 hours for the largest configuration. 

The experiments reveal that the sequential and aggregated approaches 
to configuration checking in NuSMV exhibit distinct advantages
that determine their suitability for different verification tasks.

The sequential checking method systematically explores each configuration.
Its  advantage lies in its completeness: it provides 
a comprehensive list of all reachable configurations, ensuring that no 
valid state is omitted. This  nature makes 
the sequential approach particularly valuable for analysis, where 
the generated configurations can be further examined using simulation 
to assess structural properties, other criteria. 
However, this completeness comes at a computational cost. 
Since each configuration is verified independently, 
the overall process tends to be  slower than the aggregated approach,
especially as the number of configurations increases.

In contrast, the aggregated configuration strategy check multiple 
configurations in single verification model. This integrated formulation 
enables NuSMV to handle the problem more efficiently, resulting in 
faster verification times in most cases. 
However, the aggregated model typically returns only 
the first reachable configuration that satisfies the given specification. 
Discovering additional configurations requires iteratively negating previously 
found solutions and re-running the verification process, adds procedural 
complexity. 
On the other hand, this approach supports fine-grained control 
over the structure of configurations, as it allows the inclusion of 
constraints on specific blocks within the sequence through type-defining 
variables in NuSMV. This feature makes aggregated checking well-suited 
for  analysis on design stage, where constraints  play a central role.

The choice between these two methods should therefore depend on the verification objective: 
whether the goal is to obtain a full configuration space for detailed analysis 
or to achieve rapid validation of specific structural constraints within 
a large and complex model.

\section{Evaluation of 2D reachability checking computation time}

We performed a number of verification experiments using the workflow described in 
Chapter \ref{ch:global} to validate the model and to analyze 
the robot's behavior in various conditions.

% The experiments served two primary purposes: (1) Validation of the 
% model's correctness against expected outcomes in known scenarios, 
% and (2) Performance assessment to understand how the complexity scales 
% with field size and robot length.
%
% Model Validation in Various Scenarios: We tested the model on grid sizes 
% ranging from very small (3x3, 4x4) up to moderate sizes like 7x7. 
% The robot’s assembly sequence length was also varied (we tried sequences of 
% length 2 up to 6 or more parts). We placed parts and obstacles in different 
% configurations, including edge cases:
%
% Scenarios where the robot can immediately find all parts (straightforward success).
%
% Scenarios where one of the required parts is present but located behind 
% obstacles, forcing the robot to take a longer path.
%
% Scenarios where a required part is completely absent from the field (
% which should inevitably lead to failure).
%
% Cases where the robot is boxed in or starts in a tight area, to see 
% if the model correctly identifies the failure (no path) early.
%
% Random distributions of a variety of parts on the field, with the robot 
% tasked to assemble one particular sequence out of the clutter.
%
% In each case, we checked that the counterexample path returned by nuXmv 
% made sense and corresponded to a feasible execution of the intended rules
% . The correctness of the NuSMV model’s logic was thus empirically confirmed 
% by these tests: for instance, in scenarios where a part was unreachable, 
% the model always eventually produced a FAILED trace (and the trace clearly 
% showed the robot searching all around an obstacle until no options remained
% ). In scenarios where multiple paths to a part existed, the model checker 
% inherently found one path (not necessarily all) — typically it finds one 
% of the shortest paths due to the way bounded search works in nuXmv, 
% but we did not formally ensure minimality. We cross-validated the results 
% using a manual simulation in a few cases and also with a custom visualization: 
% the GUI’s drawn path allowed us to see if the robot took a reasonable route
% . All these tests gave confidence that the model accurately implements the 
% intended behavior of the SRRS under our assumptions, and that the automation 
% correctly captures the traces.

We ensured that the automatic generation itself was correct by 
comparing an automatically generated model (for a simple scenario) 
to the hand-written initial model. They matched in structure, and running 
both yielded the same outcome. 
% This was aided by using the Formal-Methods Playground tool created by \c
% for quick iteration; we would generate a model with Python, then load it in the Playground or directly in nuXmv 
% to verify it. The Formal-Methods Playground [2] provided a convenient environment 
% to develop the initial model and even visualize state space for small 
% cases, which complemented our Python automation approach.

Using a simple Python script, we measured the 
time required for model generation, model checking, and counterexample processing 
under different scenario sizes. 
We were particularly interested in how the computation grows with field size
and with robot size.

\subsubsection{ Field Size }
We kept the robot’s sequence length fixed 
and increased the grid size (4x4, 5x5, 6x6, and 7x7)
to see how the state space explosion affects runtime. 
The results, shown in Figure \ref{fig:2DtimeField}, 
indicated that the verification time grows quickly with 
the area of the grid. 

\begin{figure}[h!]
    \centering
    \includegraphics[height=8cm]{2Dreachability/TimeFielsSize.png} 
    \caption{Model checking time with field size}
    \label{fig:2DtimeField}
\end{figure}

A larger field means more possible positions 
and more possible paths for the robot, which dramatically increases the 
state space that the model checker must explore. We found that going beyond 
a 7x7 grid with a 4-part robot made NuSmv execution considerably slower 
in our model. The relationship was not strictly exponential in 
field size for our tested range, but it was steep—significantly higher 
runtime for each increment in field dimension.

\subsubsection{Robot Length }
We then fixed a field size to 6 and varied the number of parts 
in the robot's sequence length. The computation times for 
model generation, verification, and trace processing as a function of 
robot size are plotted in Figure \ref{fig:2DtimeRobot}. 

\begin{figure}[h!]
    \centering
    \includegraphics[height=8cm]{2Dreachability/TimeRobotSize.png} 
    \caption{Model checking time with robot size}
    \label{fig:2DtimeRobot}
\end{figure}

As expected, longer sequences also 
increase the state space, because the assembly\_index has a larger range 
and more parts need to be found in order. 
This tends to increase the 
length of counterexample traces.
With a 6x6 field, going from a robot requiring 
3 parts to one requiring 5 parts made the verification slower, but the 
impact was much less pronounced than during increasing the field size. 
This suggests that spatial complexity (field area) was the more limiting 
factor in our model than the temporal length of the task.

The performance data was collected on a standard laptop with Intel Core i7 
2.4 GHz, running Ubuntu 24.04. We recorded the times for model generation,
which were negligible for these sizes, on the order of fractions of 
a second up to a couple of seconds for the largest cases, model checking,
which dominated the time, ranging from under a second for very small 
cases to many seconds or minutes for the larger ones, and counterexample processing,
also very fast, since parsing a text trace is trivial compared to 
the model checking. 

One qualitative observation from all our experiments is that the counterexample 
produced by NuSmv was often the shortest possible path to reach the 
goal or failure. That is, the model checker tends to find an optimal or 
near-optimal solution in terms of number of steps.
% , likely because any 
% unnecessary wandering would constitute a longer counterexample and model 
% checking algorithms often find counterexamples by BFS or iterative deepening 
% (thus finding shorter traces first). 
For example, if there was a more 
direct route for the robot to gather parts, the first counterexample usually 
reflected that direct route rather than a circuitous one. In our tests,
for all evaluated configurations up to length-4 robots in 5x5 grids 
(and many in 6x6), the first counterexample found corresponded to a 
strategy where the robot made efficient moves, not visiting irrelevant areas.
This suggests that, at least for these sizes, our approach could be 
leveraging the model checker as not just a verifier but also a planner that 
yields shortest paths.
However, we did not formally guarantee optimality, 
and it is possible for a model checker to output a non-shortest counterexample 
in some cases.
% (especially if heuristics or BDD search orders favor a 
% particular exploration). 
% We simply noted this behavior; 
A more rigorous analysis 
could be done to confirm path optimality or lack thereof. 
% Regardless, 
% having a concrete execution trace (shortest or not) is extremely valuable 
% for interpreting the system’s behavior.

In summary, the verification experiments confirmed that the NuSMV model 
correctly captures the reachability of both success and failure conditions 
in a self-replicating robot scenario. The automated workflow significantly reduced 
the effort to test different scenarios, and the results provided  
confidence in the model 
% and insight into performance 
% bottlenecks (through timing measurements). 
Figures \ref{fig:2DtimeField} and \ref{fig:2DtimeRobot} provide a 
high-level view of how scaling the problem size impacts the computational cost,
which is important when considering more complex or larger-scale applications 
of this verification approach.


\section{Threats to validity}

In our work, the integration of formal verification through 
NuSMV, simulation in ROS2 and Gazebo, and the conceptual modeling of 
self-replicating robotic system introduces several sources of 
uncertainty and possible threats to validity.
The purpose of this section is to critically evaluate these threats and 
to define their implications.
% in terms of internal validity -- whether the 
% results accurately represent the studied system, external validity --
% the generalizability of findings, construct validity  -- whether the models 
% represent the intended concepts, and conclusion validity --the reliability of 
% analytical results.
Understanding these potential weaknesses is essential not only for interpreting 
the results  but also for guiding future work aimed at reducing their 
impact. \\

Three main sources of validity threats were identified 
in the developed system:
\begin{enumerate}
    \item Limited generalization of the developed system — the applicability of 
    the proposed verification methods is confined to a specific class of 
    SRRS systems.
    \item Discretization errors and computational trade-offs — the accuracy and 
    performance of the model are constrained by the discretization parameter.
    \item Inconsistencies between NuSMV models, simulation environments, and real-world 
    -- abstractions and simplifications lead to mismatches between modeled 
    and physical behaviors.
\end{enumerate}

\subsubsection{ Limited generalization of the developed system }

The approaches and implementations presented in Chapter \ref{ch:implementation} and \ref{ch:simulation}
were specifically designed and validated on modular, self-replicating robotic 
systems with relative robot (robots using relative motion between their parts 
and their environment to move and build structures). 
These systems are characterized by discrete structural blocks,
that can be mechanically connected, detached, and rearranged to form new robots.
The model checking framework assumes a 
relative coordinate representation, where all positional and kinematic relations 
are expressed relative to the robot's local reference frame.

While this abstraction enabled a compact and computationally tractable formal 
model, it  limits external validity. 
In particular, the proposed verification methods may not directly generalize to
non-modular SRRS architectures, such as continuous-body or deformable robots,
self-replicating swarm systems, where replication emerges from collective behavior or
 hybrid robotic systems, where the replication process 
involves chemical or additive-manufacturing operations.

For such systems, the assumptions about structure, connectivity, and 
discrete transitions used in this work would not necessarily hold. 
Additionally, the temporal logic specifications we used depend on the ability to 
enumerate discrete states.  SRRS with other type of structure 
 would therefore require reformulation of the model semantics. 

Therefore, the findings of this work should be interpreted as representative 
of a specific class of SRRS rather than as a universal model checking methodology 
for all self-replicating robots.

\subsubsection{ Discretization errors}

The second threat to validity arises from the discretization parameter 
employed during model generation. In formal verification of  
motion-related properties, the continuous workspace of the robot must 
be discretized into a finite set of states to make model checking computationally 
feasible. 
In our system, each step of discretization defines a quantization 
of angular  and  linear space,  determining how finely 
the robot's configuration space is represented.

This choice introduces the  accuracy–efficiency contradiction:     

A coarse discretization with a large step size reduces the number of 
states, leading to faster model checking, but at the expense of precision. 
A fine discretization with a small step size increases model 
fidelity but causes exponential growth in the number of states, leading 
to  extremely long verification times. 

These discretization-related errors directly threaten the validity,
as the model's discrete representation might fail to capture the intended 
continuous behaviors. 

This problem can be mitigated by performing  multi-resolution analysis, 
where verification is performed at multiple levels of discretization and 
results are compared for consistency.
Using abstraction refinement techniques, in which coarse verification results 
guide selective refinement of critical regions in the state space.

Future research should therefore formalize the relationship between discretization,
verification precision, and computational time for SRRS,
to develop flexible model parameterization.

\subsubsection{Inconsistencies between NuSMV model, simulation, and reality}

Another threat to validity stems from inconsistencies between the 
used NuSMV model, the ROS2–Gazebo simulation, and the real-world physical implementation 
of the robot. 
While the verification model captures the logical and 
structural relations among robot components, certain geometric and kinematic 
simplifications were made to keep the model tractable.

A concrete example arises in the representation of blocks of type block\_pitch, 
where in the simulation and real robot design an axis of rotation  is 
shifted along the rotation axis relative to the connected block. 
This shift introduces a small translational displacement during rotation—
an effect that is physically significant for precise assembly but was 
omitted in the NuSMV model for simplicity. 
As a result, the formal model assumes perfectly aligned rotational joints,
whereas the physical robot exhibits a small offset.

The potential consequence is that a configuration deemed 
reachable in the formal model might not be physically realizable due 
to geometric interference or misalignment.

To mitigate this threat we used an alternating shift direction in the simulation, 
and in future this and similar aspects can be incorporated into the NuSMV model.

Performing cross-validation between model checking and simulation, 
as partially implemented in this thesis in Chapters \ref{ch:validationCounterexample}
\ref{ch:validationConfiguration},
allows detecting inconsistencies by replaying verified sequences 
in the simulation environment.

% While these mitigations can reduce the gap, a complete alignment between 
% symbolic and physical models is challenging, particularly in modular robotics 
% where even minor geometric deviations can propagate through the assembly
% . Consequently, verification results should be interpreted with an 
% awareness that they reflect idealized structural conditions.

\subsubsection{Summary}

The three identified sources of validity threats represent the main 
factors that may constrain the implementation to other sysems and interpretation of 
the results presented in this thesis. 
None of these threats invalidate the fundamental conclusions of 
the research, but they define the operational boundaries within which 
those conclusions remain valid.

From a methodological perspective, these threats highlight the necessity of 
maintaining a multi-layered validation pipeline, where formal verification, 
simulation-based testing, and physical experimentation complement one 
another. 
Model checking provides exhaustive logical guarantees within its 
abstract domain, while simulation validates these results under physical 
dynamics, and real-world testing confirms them under environmental uncertainty.
By addressing these aspects, subsequent studies can move toward fully 
validated, scalable frameworks for self-replicating robotic systems.
